\documentclass[11pt,a4paper]{article}
% Shared preamble for RuPhO English problem statements (match T1 2026 style).
% Use: \input{latex/rupho_en_preamble.tex} after \documentclass.
\usepackage[margin=1in]{geometry}
\usepackage{amsmath,amssymb}
\usepackage{graphicx}
\usepackage{float}
\usepackage{enumitem}
\usepackage{microtype}
\usepackage{caption}
\usepackage{tabularx}
\usepackage{hyperref}

\captionsetup{font=small,labelfont=bf,skip=6pt}

\setlist[itemize]{leftmargin=1.2cm}
\setlist[enumerate]{leftmargin=1.2cm}

% One outer box per subproblem: [id | statement | points] (no inner rules)
\newcommand{\problembox}[3]{%
  \par\addvspace{0.42em}%
  \noindent
  \setlength{\fboxsep}{7.5pt}%
  \setlength{\fboxrule}{0.95pt}%
  \fbox{%
    \begin{minipage}{\dimexpr\linewidth-2\fboxsep-2\fboxrule\relax}
    \begin{tabularx}{\linewidth}{@{}>{\raggedright\arraybackslash\bfseries}p{2.1em} @{\hspace{0.9em}} >{\raggedright\arraybackslash}X @{\hspace{0.9em}} >{\raggedleft\arraybackslash\bfseries}p{2.4em}@{}}
    #1 & #2 & #3
    \end{tabularx}
    \end{minipage}%
  }%
  \par\addvspace{0.32em}%
}


\title{\textbf{Particles on a straight line and plane}}
\author{\normalsize X20 (2020) --- English translation}
\date{}
\begin{document}
\maketitle

There is a straight line $l$ and a plane $a$. They are perpendicular. Particle 1 can move in a straight line, and particle 2 can move along a plane; the mass of each particle is $m$. Particles are attracted according to the law of gravity. There is no friction.

\begin{figure}[H]
    \centering
    \includegraphics[width=0.72\textwidth]{T4_EN_assets/fig1.jpg}
    \caption{Figure 1.}
\end{figure}

\problembox{A1}{Initially, particle 1 is at rest, and the second is set in motion. Its speed is perpendicular to the line connecting the particles, and the distance between the particles is $R$. Find the minimum value of the speed of the second particle $v_{02}$, at which the particles will be at an infinite distance from each other.}{3.00}

Now the particles are at equal distances from the point of intersection of the line and the plane, the distance between them is $R$. The particles are given equal velocities $v_0$, lying in the same plane perpendicular to the plane $\alpha$.

\problembox{A2}{Find the value $v_0$ at which particle 2, after a long time, will finally stop at a finite distance from the origin and the value of the speed $v_1$ of the first particle in this case.}{3.00}

\vspace{1em}
\noindent\textit{Source:} \url{https://pho.rs/p/15}

\end{document}